\documentclass[a4paper,fontsize=11pt,fleqn]{scrartcl}

\usepackage{../../../Styles/Eval}

\newcommand{\chnum}{Chapitre 5 et 8}
\newcommand{\chtitle}{Corrigé DS5 -- Polarité des molécules et titrages colorimétriques}
\newcommand{\dtype}{Corrigé}

\begin{document}
\maketitle

\section{Polarité des molécules organiques}

On considère les espèces chimiques suivantes : \ce{HSO4^-}, \ce{HPO3} et \ce{H2CO2}.

\begin{enumerate}
    \item Proposer pour chaque molécule un schéma de Lewis correct.
    \item En se basant sur le tableau des électronégativités de Pauling, réécrire les structures de Lewis en indiquant les éventuelles charges partielles qui y sont présentes.
    \item Déterminer la géométrie de ces structures autour des atomes de soufre, de phosphore et du carbone.
    \item Déterminer si chacune des molécules est polaire ou apolaire.
\end{enumerate}

\subsection*{Corrigé}

\begin{enumerate}
    \item \textbf{Schémas de Lewis}

    \begin{itemize}
        \item \textbf{\ce{HSO4^-}} : l'atome central est le soufre.  
        Une représentation correcte est :
        \[
        \ce{HO-S(=O)2-O^-}
        \]
        Le soufre est lié à 4 oxygènes : un groupe \ce{-OH}, deux oxygènes en double liaison et un oxygène portant la charge négative.

        \item \textbf{\ce{HPO3}} : dans le cadre de l'exercice, on peut écrire :
        \[
        \ce{HO-P(=O)2}
        \]
        Le phosphore est lié à trois oxygènes, dont un sous forme \ce{-OH}.

        \item \textbf{\ce{H2CO2}} : on reconnaît l'acide méthanoïque :
        \[
        \ce{H-C(=O)-OH}
        \]
        Le carbone est lié à un hydrogène, à un oxygène par double liaison et à un groupe \ce{-OH}.
    \end{itemize}

    \item \textbf{Charges partielles}

    Les liaisons \ce{O-H}, \ce{S-O}, \ce{P-O} et \ce{C-O} sont polarisées car l'oxygène est plus électronégatif que H, S, P et C.

    \begin{itemize}
        \item Dans toutes les liaisons avec l'oxygène, l'oxygène porte une charge partielle \(\delta^-\).
        \item Les atomes liés à l'oxygène (H, S, P, C) portent une charge partielle \(\delta^+\).
    \end{itemize}

    On peut donc indiquer :
    \begin{itemize}
        \item pour \ce{HSO4^-} : \(\ce{S^{\delta+}}\) et les oxygènes \(\ce{O^{\delta-}}\), ainsi que l'hydrogène du groupe \ce{OH} en \(\delta^+\) ;
        \item pour \ce{HPO3} : \(\ce{P^{\delta+}}\), oxygènes en \(\delta^-\), hydrogène du groupe \ce{OH} en \(\delta^+\) ;
        \item pour \ce{H2CO2} : \(\ce{C^{\delta+}}\), oxygènes en \(\delta^-\), hydrogène du groupe \ce{OH} en \(\delta^+\).
    \end{itemize}

    \item \textbf{Géométrie autour de l'atome central}

    \begin{itemize}
        \item \textbf{\ce{HSO4^-}} : autour du soufre, il y a 4 domaines électroniques liants.  
        Géométrie : \textbf{tétraédrique}.

        \item \textbf{\ce{HPO3}} : autour du phosphore, il y a 3 domaines électroniques liants.  
        Géométrie : \textbf{trigonale plane}.

        \item \textbf{\ce{H2CO2}} : autour du carbone, il y a 3 domaines électroniques liants.  
        Géométrie : \textbf{trigonale plane}.
    \end{itemize}

    \item \textbf{Polarité}

    \begin{itemize}
        \item \ce{HSO4^-} est \textbf{polaire} ;
        \item \ce{HPO3} est \textbf{polaire} ;
        \item \ce{H2CO2} est \textbf{polaire}.
    \end{itemize}

    En effet, ces espèces comportent des liaisons polarisées et leur géométrie ne permet pas une compensation totale des moments dipolaires.
\end{enumerate}

\newpage

\section{Dosage de l'eau de Javel}

L'eau de Javel est une solution oxydante utilisée comme désinfectant et comme décolorant. Selon l'emploi désiré, il est important de connaître sa concentration.

\begin{documentboxenv}{Synthèse de l'eau de Javel}
Charles Guillaume Scheele (1742-1786), pharmacien suédois, découvrit le dichlore au 18\textsuperscript{e} siècle. Ce gaz intervient dans la fabrication de l'eau de Javel, qui doit son nom à un ancien village aujourd'hui intégré à un quartier de Paris.\\
C'est à Javel que Claude Louis Berthollet (1748-1822), directeur à la manufacture des Gobelins, fabriqua ce produit décolorant et désinfectant et l'employa en 1785 pour le blanchiment des toiles textiles. L'eau de Javel est obtenue par action du dichlore \ce{Cl2(g)} sur l'hydroxyde de sodium et contient notamment des ions \ce{Cl^-(aq)}, \ce{ClO^-(aq)}, \ce{Na^+(aq)} ainsi que de l'eau \ce{H2O(l)}.\\
En milieu acide, l'eau de Javel subit une transformation complète représentée par :
\[
\ce{ClO^-(aq) + Cl^-(aq) + 2H^+(aq) -> H2O(l) + Cl2(g)}
\]
Cette transformation permet de définir le degré chlorométrique. Celui-ci est égal au volume, exprimé en litres, de dichlore produit par un litre d'eau de Javel. Ce volume est mesuré à \SI{0}{\celsius} sous une pression de \SI{1.013}{\bar}.
\end{documentboxenv}

\begin{documentboxenv}{Méthode expérimentale}
Pour vérifier l'indication portée sur une bouteille commerciale d'eau de Javel, \SI{12}{\degree chl}, on procède à une expérience.\\
\textbf{Principe de la manipulation}
\begin{itemize}
    \item On ajoute un excès d'ions iodure à un volume connu de solution d'eau de Javel.
    \item En milieu acide, les ions hypochlorite \ce{ClO^-(aq)} oxydent les ions iodure \ce{I^-(aq)}.
    \item L'équation de la réaction est :
    \[
    \ce{ClO^-(aq) + 2I^-(aq) + 2H^+(aq) -> I2(aq) + H2O(l)}
    \]
\end{itemize}
On considère cette réaction comme totale. Le diiode formé est ensuite dosé par les ions thiosulfate.
\end{documentboxenv}

\subsection{Réaction entre les ions hypochlorite et les ions iodure}

\begin{enumerate}
    \item L'eau de Javel commerciale étant trop concentrée, il faut d'abord effectuer une dilution du dixième pour obtenir \SI{50,0}{mL} de solution diluée $S$. Décrire une méthode qui permet d'effectuer cette dilution. Préciser la verrerie nécessaire (noms et volumes).
    \item Dans un erlenmeyer (noté $E$), on introduit dans cet ordre : $V = \SI{10,0}{mL}$ de solution $S$ puis $V' = \SI{20,0}{mL}$ de solution d'iodure de potassium. Indiquer la verrerie à utiliser.
\end{enumerate}

\subsection*{Corrigé}

\begin{enumerate}
    \item \textbf{Dilution au dixième}

    Pour préparer \SI{50,0}{mL} de solution diluée au dixième :
    \begin{itemize}
        \item on prélève \SI{5,0}{mL} de solution commerciale avec une \textbf{pipette jaugée de \SI{5,0}{mL}} ;
        \item on verse ce prélèvement dans une \textbf{fiole jaugée de \SI{50,0}{mL}} ;
        \item on complète avec de l'eau distillée jusqu'au trait de jauge ;
        \item on bouche et on homogénéise.
    \end{itemize}

    \item \textbf{Verrerie à utiliser}

    \begin{itemize}
        \item pour prélever \SI{10,0}{mL} de solution $S$ : \textbf{pipette jaugée de \SI{10,0}{mL}} ;
        \item pour prélever \SI{20,0}{mL} de solution d'iodure de potassium : \textbf{pipette jaugée de \SI{20,0}{mL}} ;
        \item le mélange est réalisé dans un \textbf{erlenmeyer}.
    \end{itemize}
\end{enumerate}

\subsection{Détermination de la concentration en eau de Javel}

\begin{documentboxenv}{Montage pour dosage}
On dose le diiode formé par une solution de thiosulfate de sodium de concentration molaire apportée
\[
c_1 = \SI{0,10}{\mole\per\liter}
\]
Le volume à l'équivalence est :
\[
V_{1E} = \SI{10,0}{mL}
\]
\end{documentboxenv}

\begin{enumerate}
    \item Écrire l'équation de la réaction entre le diiode et les ions thiosulfate.
    \item Justifier la décoloration progressive du mélange réactionnel.
    \item Quel est le réactif limitant avant l'équivalence ? Après l'équivalence ? À l'équivalence ?
    \item Déduire des résultats de l'expérience la quantité de matière de diiode présente dans le mélange réactionnel.
    \item Calculer la quantité de matière d'ions hypochlorite initialement présents dans le prélèvement de volume $V$.
    \item Déterminer la concentration en ions hypochlorite de la solution $S$, puis de la solution commerciale.
    \item En utilisant l'équation de la réaction chimique donnée au Doc. 1, calculer la quantité de matière de dichlore produite par un litre d'eau de Javel.
    \item En déduire le degré chlorométrique de l'eau de Javel commerciale utilisée.
\end{enumerate}

\subsection*{Corrigé}

\begin{enumerate}
    \item \textbf{Équation de la réaction de dosage}

    Le diiode réagit avec les ions thiosulfate selon :
    \[
    \ce{I2(aq) + 2S2O3^{2-}(aq) -> 2I^-(aq) + S4O6^{2-}(aq)}
    \]

    \item \textbf{Décoloration progressive}

    Le diiode en solution donne, en présence d'amidon, une coloration bleue.  
    Au cours du dosage, les ions thiosulfate consomment progressivement le diiode.  
    La quantité de diiode diminue donc, et la coloration bleue s'atténue puis disparaît à l'équivalence.

    \item \textbf{Réactif limitant}

    \begin{itemize}
        \item avant l'équivalence : les ions thiosulfate sont le \textbf{réactif limitant} ;
        \item à l'équivalence : les réactifs sont introduits dans les proportions stoechiométriques ;
        \item après l'équivalence : le diiode a été totalement consommé, c'est donc le \textbf{diiode} qui était limitant.
    \end{itemize}

    \item \textbf{Quantité de matière de diiode}

    À l'équivalence :
    \[
    n(\ce{S2O3^{2-}}) = c_1 \times V_{1E}
    \]
    avec
    \[
    c_1 = \SI{0,10}{\mole\per\liter}
    \quad \text{et} \quad
    V_{1E} = \SI{10,0}{mL} = \SI{10,0e-3}{L}
    \]
    donc
    \[
    n(\ce{S2O3^{2-}}) = 0,10 \times 10,0\times10^{-3} = \SI{1,0e-3}{mol}
    \]

    Or d'après l'équation :
    \[
    \ce{I2 + 2S2O3^{2-} -> 2I^- + S4O6^{2-}}
    \]
    donc
    \[
    n(\ce{I2}) = \frac{1}{2}n(\ce{S2O3^{2-}})
    \]
    soit
    \[
    n(\ce{I2}) = \frac{1,0\times10^{-3}}{2} = \SI{5,0e-4}{mol}
    \]

    \item \textbf{Quantité de matière d'ions hypochlorite dans le prélèvement}

    La réaction entre les ions hypochlorite et les ions iodure est :
    \[
    \ce{ClO^- + 2I^- + 2H^+ -> I2 + H2O}
    \]

    Les coefficients stoechiométriques montrent que :
    \[
    n(\ce{ClO^-}) = n(\ce{I2})
    \]
    donc dans le prélèvement de volume \(\SI{10,0}{mL}\) :
    \[
    n(\ce{ClO^-}) = \SI{5,0e-4}{mol}
    \]

    \item \textbf{Concentration en ions hypochlorite}

    Dans la solution diluée $S$ :
    \[
    c_S = \frac{n(\ce{ClO^-})}{V}
    = \frac{5,0\times10^{-4}}{10,0\times10^{-3}}
    = \SI{5,0e-2}{\mole\per\liter}
    \]

    Donc :
    \[
    c_S = \SI{5,0e-2}{\mole\per\liter}
    \]

    La solution commerciale est 10 fois plus concentrée :
    \[
    c_{\text{commerciale}} = 10 \times c_S = \SI{5,0e-1}{\mole\per\liter}
    \]

    Donc :
    \[
    c_{\text{commerciale}} = \SI{0,50}{\mole\per\liter}
    \]

    \item \textbf{Quantité de matière de dichlore produite par 1 litre d'eau de Javel}

    D'après la réaction :
    \[
    \ce{ClO^- + Cl^- + 2H^+ -> H2O + Cl2}
    \]

    Il y a une proportion \(1:1\) entre \(\ce{ClO^-}\) et \(\ce{Cl2}\).

    Dans \SI{1,0}{L} d'eau de Javel commerciale :
    \[
    n(\ce{ClO^-}) = c_{\text{commerciale}} \times V = 0,50 \times 1,0 = \SI{0,50}{mol}
    \]

    Donc :
    \[
    n(\ce{Cl2}) = \SI{0,50}{mol}
    \]

    \item \textbf{Degré chlorométrique}

    Avec le volume molaire :
    \[
    V_m = \SI{22,4}{\liter\per\mole}
    \]

    le volume de dichlore produit par \SI{1,0}{L} d'eau de Javel vaut :
    \[
    V(\ce{Cl2}) = n(\ce{Cl2}) \times V_m
    = 0,50 \times 22,4
    = \SI{11,2}{L}
    \]

    Le degré chlorométrique de l'eau de Javel est donc :
    \[
    \boxed{\SI{11,2}{\degree chl}}
    \]

    Cette valeur est proche de l'indication de \(\SI{12}{\degree chl}\).
\end{enumerate}

\end{document}